\section{Threats to Validity}
\label{sec:threats}

\paragraph{Construct validity.}
The rubric asks whether a review addresses a topic, not whether what it says
is true. A hallucinated review scored $9.2$ against a real review's $9.0$.
All rubric numbers are therefore coverage measurements; correctness claims
rest on Experiment~2's oracle. The human-study interface verifies only that a
referenced file exists, not that the change genuinely affects it.
Binary scoring and seven immovable criteria (at $0\,\%$ or $100\,\%$ in both
arms) further dilute the aggregate, which is one reason the pre-specified
subscale is reported alongside the total. The graph arm also writes the
longest reviews ($+12\,\%$ in characters; Table~\ref{tab:review-length}).
Retrieval adds far more input, writes reviews of baseline length, and still
gains most on the total, so length alone does not explain the pattern; a
length contribution to the graph-arm estimate cannot be ruled out.

\paragraph{Internal validity.}
Language models judge language-model output. The panel and the generator share
training data, so a judge may reward its own family's output. Mitigations:
the five-judge panel draws from five providers, per-judge verdicts are
published for re-aggregation, and judging runs at $T = 0.0$. Tested alone,
four of five judges reproduce the targeted effect at $p < 0.05$ and all five
preserve its direction.

The arms differ not only in the context block but also in the instruction
paragraph that names it (Appendix~\ref{app:prompts}), so the contrast
confounds content with instruction. Experiment~2's oracle mitigates this:
detection requires naming a genuinely broken site, which the instruction
alone cannot supply. The retrieval index covers only the first 200
candidate source files per repository (Section~\ref{sec:approach-rag}), so
in large repositories the retrieval arm is bounded by its index as well as
by its method.

Generation is not bit-reproducible: regenerating the graph arm from identical
inputs moves the nine-criterion score by $0.88$~points on average, unbiased
but of the same magnitude as the effect. The reviews and per-judge verdicts
are retained, so every number recomputes from stored outputs. Both
experiments generate at $T=0.0$, but they use different samples, outcomes,
and diff-length caps. Their effect magnitudes are therefore not compared
directly.

Experiment~2's injections are drawn from mutation-testing operators; selection
was direction-blind, but a seeded rename is not a defect a human introduced.
Its two-stage adjudication was introduced after an audit exposed false
negatives in the original judge prompt. The correction is post-hoc; the
frozen reviews and original verdicts remain available, and the thesis reports
the corrected measurement as such.

The human study uses an assisted interface with summaries, file tooltips, and
a diff shown by default. Its primary result therefore measures preference
under that interface, not unaided review reading. Four pre-registered
diagnostic checks of aid use, rationale content, file-count association, and
task time were not implemented and support no thesis claim.

\paragraph{External validity.}
Three substitute generators preserve the direction but none reaches
significance; both the generator and the held-out rows are from the earlier
text-heuristic configuration at $n = 40$ under a three-judge panel, so they
are directional evidence rather than tests of the current configuration. The
rubric-measured benefit is conditional on the generator.

Two graph constructions appear (the Joern code property graph in
Experiments~1 and~2, and import-statement edges in the human study), and
they are not interchangeable. The human study's edges are file-level; its
effect is an effect of that construction. Claims about graph
constructions are claims about these constructions.

All 35 pull requests were merged (survivorship). The selection acts
identically on all arms, which protects the within-sample comparison, but it
can still limit how the treatment effect generalises beyond the selected
sample. The human study covers six pull requests from three Python
libraries, so its preference result does not establish cross-language
generality.

\paragraph{Conclusion validity.}
With $n = 35$ the study is powered for $d_z \geq 0.47$ at 80\,\%. The graph arm
reaches significance on the full rubric ($p = 0.017$,
$d_z = +0.44$) but does not survive Bonferroni correction
($p_{\text{corr}} = 0.066$). The conclusion is therefore framed in
terms of the subscale, where the effect is larger
($d_z = +0.56$, $p = 0.003$,
$p_{\text{corr}} = 0.011$). The human study's 27
raters are sufficient for the primary endpoint (rank-biserial $r = 0.80$)
but explicitly underpowered for the overall-usefulness endpoint, which is
designated estimation-only. On 12 further pull requests not used in the
main analysis, evaluated under the earlier text-heuristic configuration, the
KG-relevant estimate remains positive ($+0.42$) but does not confirm the
effect ($p = 0.312$); the small replication was underpowered and is
interpreted as directional consistency rather than confirmation.
